% Source: https://artofproblemsolving.com/community/c3h1771314_donut_problem_kinda
% Stable identity: aops-topic-1771314-q01
The answer is $\boxed{\textbf{(E)}\ 7}$. (Regions are assumed connected; regions meeting only at a point may share a colour.)

\textbf{Seven colours are sometimes needed.} Tile the plane with regular hexagons centred at the points $a\mathbf u+b\mathbf v$ ($a,b$ integers) of a triangular lattice, and give the hexagon at $a\mathbf u+b\mathbf v$ the label $a+3b \bmod 7$. The six neighbours of a hexagon are at $\pm\mathbf u$, $\pm\mathbf v$, $\pm(\mathbf u-\mathbf v)$, so their labels differ from its own by $\pm1$, $\pm3$, $\pm2$: they carry all six other labels.

\begin{asy}
usepackage("amsmath");
size(6cm);
defaultpen(fontsize(10pt));
real s3=sqrt(3);
pair u=(s3,0), v=(s3/2,1.5);
path hexa(pair c){ return c+dir(90)--c+dir(150)--c+dir(210)--c+dir(270)--c+dir(330)--c+dir(30)--cycle; }
pen[] col={rgb(1,0.85,0.85),rgb(1,0.93,0.75),rgb(1,1,0.75),rgb(0.82,0.95,0.82),rgb(0.8,0.92,1),rgb(0.88,0.85,1),rgb(0.95,0.85,0.95)};
for(int a=-2;a<=2;++a) for(int bb=-2;bb<=2;++bb){
  if(abs(a+bb)>2) continue;
  pair c=a*u+bb*v;
  int k=((a+3*bb)%7+7)%7;
  filldraw(hexa(c),col[k],black+linewidth(0.7));
  label("$"+string(k)+"$",c);
}
label("\mbox{}",(min(currentpicture,user=true)+max(currentpicture,user=true))/2,black); // keeps the text after the figure black when LaTeX typesets the labels inline
\end{asy}

The labelling is unchanged by the translations $a\mathbf u+b\mathbf v$ with $a+3b\equiv0\pmod7$, for example $7\mathbf u$ and $-3\mathbf u+\mathbf v$; these form a lattice, and identifying points of the plane that differ by such a translation turns the plane into a torus (a parallelogram with opposite sides glued). On this torus there are exactly $7$ hexagons, one with each label, and each touches the other six along an edge. So this map needs $7$ colours.

\textbf{Seven colours always suffice.} Put a vertex in each region and join two vertices by an edge, drawn across their common boundary, when the regions are adjacent. This gives a simple graph $G$ drawn on the torus without crossings, and the same is true for every subgraph of $G$. For a graph with $V\ge3$ vertices, $E$ edges and $F$ faces drawn on the torus, Euler's formula for the torus gives $V-E+F\ge0$, and every face is bounded by at least $3$ edges, each edge lying on at most $2$ faces, so $3F\le2E$. Hence $E\le3V$, and the average degree $\dfrac{2E}{V}$ is at most $6$: some vertex has degree at most $6$.

Now colour by induction on the number of vertices. Remove a vertex $x$ of degree at most $6$, colour the rest with $7$ colours (the smaller graph is again drawn on the torus), and give $x$ a colour not used by its at most $6$ neighbours. Graphs with at most $7$ vertices are trivially $7$-colourable. Hence every map on the torus can be coloured with $7$ colours, and $7$ is the minimum. $\blacksquare$
